%% ======================================================================
%% TCC - Camila Sabino Ribeiro
%% Formatação conforme o Manual para Elaboração de Trabalhos Científicos
%% do Ibmec (4ª ed., 2022), baseado na ABNT NBR 14724.
%%
%% Compilar com pdfLaTeX + Biber (biblatex-abnt). A fonte Arial é
%% substituída por Helvetica (Nimbus Sans).
%% ======================================================================

% twoside + openright: frente e verso, capítulos começam em página ímpar
% (o manual recomenda). Se for entregar só em PDF/uma face, troque por:
% \documentclass[12pt,a4paper,oneside,openany]{report}
% e, no geometry, use "left=3cm,right=2cm" no lugar de inner/outer.
\documentclass[12pt,a4paper,twoside,openright]{report}

\usepackage[T1]{fontenc}
\usepackage[utf8]{inputenc}
\usepackage[english,brazilian]{babel}
\usepackage{helvet}                       % Arial-like
\renewcommand{\familydefault}{\sfdefault}
\usepackage{microtype}

% Margens: anverso esquerda/superior 3 cm, direita/inferior 2 cm;
% verso direita/superior 3 cm, esquerda/inferior 2 cm.
\usepackage[a4paper,twoside,inner=3cm,outer=2cm,top=3cm,bottom=2cm,
            headheight=14pt,headsep=0.5cm]{geometry}

\usepackage{setspace}
\usepackage{indentfirst}
\usepackage{graphicx}
\usepackage{placeins}
\usepackage{booktabs}
\usepackage{amsmath}
\usepackage{amssymb}
\usepackage{csquotes}
\usepackage{import}
\usepackage{fancyhdr}
\usepackage[hang,flushmargin]{footmisc}
\usepackage{titlesec}
\usepackage[font=footnotesize,labelsep=endash,position=top]{caption}
\usepackage{newfloat}
\usepackage{titletoc}
\usepackage{float}
\usepackage{longtable}
\usepackage{array}
\usepackage{caption}

% Citações e referências no padrão ABNT (biblatex-abnt + Biber).
% \cite{chave}       -> (AKGIRAY, 1989)
% \citeonline{chave} -> Akgiray (1989)   (se der erro, use \textcite{chave})
\usepackage[backend=biber,style=abnt]{biblatex}
\addbibresource{referencias.bib}

\usepackage{hyperref}
\hypersetup{
    colorlinks=true,
    linkcolor=black,
    citecolor=black,
    urlcolor=black,          % manual: texto em cor preta
    plainpages=false,
    pdftitle={TCC Camila},
    pdfauthor={Camila Sabino Ribeiro}
}

%% ----------------------------------------------------------------------
%% DADOS DO TRABALHO (edite aqui)
%% ----------------------------------------------------------------------
\newcommand{\tccinstituicao}{IBMEC BRASÍLIA}
\newcommand{\tcccurso}{CURSO DE CIÊNCIAS ECONÔMICAS}
\newcommand{\tccautor}{CAMILA SABINO RIBEIRO}
\newcommand{\tcctitulo}{Abordagens model-based e data-driven na previsão da volatilidade cambial: evidências para a taxa USD/BRL}
\newcommand{\tccsubtitulo}{}            % se houver; sai após dois pontos
\newcommand{\tcclocal}{BRASÍLIA}
\newcommand{\tccano}{2026}
\newcommand{\tccorientador}{Prof. [Esp./Ms./Dr.] Laerte Jun Takeuti}
\newcommand{\tccnatureza}{Trabalho de Conclusão de Curso apresentado ao curso
de Ciências Econômicas do Ibmec Brasília, como requisito parcial para a
obtenção do grau de Bacharel em Ciências Econômicas.}

\newcommand{\tcctitulocompleto}{\tcctitulo\ifx\tccsubtitulo\empty\else: \tccsubtitulo\fi}

%% ----------------------------------------------------------------------
%% FORMATAÇÃO GERAL
%% ----------------------------------------------------------------------
\onehalfspacing                              % entrelinhas 1,5
\setlength{\parindent}{1.25cm}               % recuo de parágrafo
\setlength{\parskip}{0pt}
\setcounter{secnumdepth}{4}                  % numeração até a seção quinária
\setcounter{tocdepth}{4}

% Nomes em português
\addto\captionsbrazilian{%
  \renewcommand{\listfigurename}{Lista de figuras}%
  \renewcommand{\listtablename}{Lista de tabelas}%
}

% Quadros e gráficos como ilustrações com lista própria
\DeclareFloatingEnvironment[fileext=loq,listname={Lista de quadros},name=Quadro]{quadro}
\DeclareFloatingEnvironment[fileext=log,listname={Lista de gráficos},name=Gráfico]{grafico}

% Fonte (obrigatória abaixo de figuras, quadros, gráficos e tabelas,
% mesmo quando for elaboração própria). Uso: \fonte{Elaborada pela autora.}
\newcommand{\fonte}[1]{\par\vspace{2pt}{\centering\footnotesize Fonte: #1\par}}

% Citação direta longa (mais de 3 linhas): recuo 4 cm, fonte 10,
% espaço simples, sem aspas.
\newenvironment{citacaolonga}
  {\par\vspace{0.5\baselineskip}\begin{singlespace}\footnotesize
   \list{}{\setlength{\leftmargin}{4cm}\setlength{\rightmargin}{0pt}%
           \setlength{\topsep}{0pt}\setlength{\listparindent}{0pt}%
           \setlength{\itemindent}{0pt}}\item\relax}
  {\endlist\end{singlespace}\par\vspace{0.5\baselineskip}}

% Nota de rodapé: fonte 10 (footnotesize em 12pt) e filete de 5 cm
\renewcommand{\footnoterule}{\vspace*{-3pt}\noindent\rule{5cm}{0.4pt}\vspace*{2.6pt}}

%% ----------------------------------------------------------------------
%% TÍTULOS DE SEÇÃO (NBR 6024)
%%  primária:    negrito, CAIXA ALTA
%%  secundária:  sem negrito, CAIXA ALTA
%%  terciária:   negrito, caixa baixa
%%  quaternária: itálico, caixa baixa
%%  quinária:    normal, caixa baixa
%%  sem número (resumo, sumário, referências...): centralizado, negrito, caixa alta
%% ----------------------------------------------------------------------
\titleformat{\chapter}[hang]{\normalfont\normalsize\bfseries}{\thechapter}{0.5em}{\MakeUppercase}
\titleformat{name=\chapter,numberless}[block]{\normalfont\normalsize\bfseries\centering}{}{0pt}{\MakeUppercase}
\titleformat{\section}[hang]{\normalfont\normalsize}{\thesection}{0.5em}{\MakeUppercase}
\titleformat{\subsection}[hang]{\normalfont\normalsize\bfseries}{\thesubsection}{0.5em}{}
\titleformat{\subsubsection}[hang]{\normalfont\normalsize\itshape}{\thesubsubsection}{0.5em}{}
\titleformat{\paragraph}[hang]{\normalfont\normalsize}{\theparagraph}{0.5em}{}

\titlespacing*{\chapter}{0pt}{0pt}{1.5\baselineskip}
\titlespacing*{\section}{0pt}{1.5\baselineskip}{1.5\baselineskip}
\titlespacing*{\subsection}{0pt}{1.5\baselineskip}{1.5\baselineskip}
\titlespacing*{\subsubsection}{0pt}{1.5\baselineskip}{1.5\baselineskip}
\titlespacing*{\paragraph}{0pt}{1.5\baselineskip}{1.5\baselineskip}

%% Sumário com a mesma tipografia dos títulos
\titlecontents{chapter}[1.8em]{\bfseries\addvspace{6pt}}
  {\contentslabel{1.8em}\MakeUppercase}{\hspace*{-1.8em}\MakeUppercase}
  {\titlerule*[0.7pc]{.}\contentspage}
\titlecontents{section}[2.8em]{}
  {\contentslabel{2.8em}\MakeUppercase}{\hspace*{-2.8em}\MakeUppercase}
  {\titlerule*[0.7pc]{.}\contentspage}
\titlecontents{subsection}[3.8em]{\bfseries}
  {\contentslabel{3.8em}}{\hspace*{-3.8em}}
  {\titlerule*[0.7pc]{.}\contentspage}
\titlecontents{subsubsection}[4.9em]{\itshape}
  {\contentslabel{4.9em}}{\hspace*{-4.9em}}
  {\titlerule*[0.7pc]{.}\contentspage}
\titlecontents{paragraph}[6em]{}
  {\contentslabel{6em}}{\hspace*{-6em}}
  {\titlerule*[0.7pc]{.}\contentspage}

%% ----------------------------------------------------------------------
%% PAGINAÇÃO
%% Todas as folhas contam a partir da folha de rosto, mas só aparece o número
%% a partir da Introdução, no canto superior direito (anverso) ou esquerdo
%% (verso), fonte 10.
%% ----------------------------------------------------------------------
\pagestyle{plain}
\fancypagestyle{plain}{\fancyhf{}\renewcommand{\headrulewidth}{0pt}\renewcommand{\footrulewidth}{0pt}}

\newcommand{\paginacaonumerada}{%
  \fancypagestyle{plain}{%
    \fancyhf{}%
    \fancyhead[RO,LE]{\footnotesize\thepage}%
    \renewcommand{\headrulewidth}{0pt}\renewcommand{\footrulewidth}{0pt}}%
  \pagestyle{plain}}

% Páginas em branco inseridas pelo openright ficam sem número
\makeatletter
\def\cleardoublepage{\clearpage\if@twoside\ifodd\c@page\else
  \thispagestyle{empty}\hbox{}\newpage\fi\fi}
\makeatother

%% Apêndices no formato "APÊNDICE A – TÍTULO" (também no sumário)
\newcommand{\iniciarapendices}{%
  \cleardoublepage
  \appendix
  \titleformat{\chapter}[hang]{\normalfont\normalsize\bfseries}
    {APÊNDICE \thechapter\ \textendash\ }{0pt}{\MakeUppercase}%
  \titlecontents{chapter}[0pt]{\bfseries\addvspace{6pt}}
    {APÊNDICE \thecontentslabel\ \textendash\ \MakeUppercase}{\MakeUppercase}
    {\titlerule*[0.7pc]{.}\contentspage}}

\begin{document}

%% ======================================================================
%% ELEMENTOS PRÉ-TEXTUAIS
%% ======================================================================

%% --- CAPA (não conta página) -----------------------------------------
\begin{titlepage}
  \centering\bfseries\singlespacing
  {\tccinstituicao\par}
  \vspace{0.5\baselineskip}
  {\tcccurso\par}
  \vspace{3cm}
  {\tccautor\par}
  \vfill
  {\tcctitulocompleto\par}
  \vfill
  {\tcclocal\par}
  {\tccano\par}
\end{titlepage}

%% --- FOLHA DE ROSTO (página 1) ---------------------------------------
\thispagestyle{empty}
\begin{center}\bfseries\singlespacing
  {\tccautor\par}
  \vspace{4cm}
  {\tcctitulocompleto\par}
\end{center}
\vspace{2cm}
\hfill\begin{minipage}{0.5\textwidth}% do meio da página para a margem direita
  \singlespacing\raggedright\normalfont
  \tccnatureza\par
  \vspace{\baselineskip}
  Orientador: \tccorientador.
\end{minipage}
\vfill
\begin{center}\bfseries\singlespacing
  {\tcclocal\par}
  {\tccano\par}
\end{center}
\newpage

%% --- VERSO DA FOLHA DE ROSTO: FICHA CATALOGRÁFICA --------------------
% Substitua pela ficha gerada pela biblioteca (fonte 10, espaço simples).
\thispagestyle{empty}
\vspace*{\fill}
\begin{center}
\fbox{\begin{minipage}{12.5cm}\footnotesize\singlespacing\raggedright
  [Inserir aqui a ficha catalográfica fornecida pela Biblioteca]
\end{minipage}}
\end{center}
\newpage

%% --- FOLHA DE APROVAÇÃO ----------------------------------------------
\newcommand{\assinatura}[2]{%
  \vspace{1.2cm}\par
  \noindent\hfill\begin{minipage}{0.8\textwidth}\centering\singlespacing
  \rule{\linewidth}{0.4pt}\par #1\par #2\end{minipage}\hfill\null\par}

\thispagestyle{empty}
\begin{center}\bfseries\singlespacing
  {\tccautor\par}
  \vspace{2cm}
  {\tcctitulocompleto\par}
\end{center}
\vspace{1.5cm}
\hfill\begin{minipage}{0.5\textwidth}
  \singlespacing\raggedright\normalfont
  \tccnatureza\par
  \vspace{\baselineskip}
  Orientador: \tccorientador.
\end{minipage}
\vspace{1.5cm}

\noindent Aprovada em: \underline{\hspace{0.7cm}}/\underline{\hspace{0.7cm}}/\underline{\hspace{1.2cm}}

\vspace{1cm}
\begin{center}\bfseries BANCA EXAMINADORA\end{center}
\assinatura{\tccorientador\ (Orientador)}{Ibmec Brasília}
\assinatura{Prof. [titulação] [Nome do examinador 1]}{[Instituição]}
\assinatura{Prof. [titulação] [Nome do examinador 2]}{[Instituição]}
\newpage

%% --- DEDICATÓRIA / AGRADECIMENTOS / EPÍGRAFE (opcionais) --------------
% \chapter*{Agradecimentos}
% Texto dos agradecimentos.
% \newpage

%% --- RESUMO (150 a 500 palavras, parágrafo único, espaço simples) ----
\cleardoublepage
\chapter*{Resumo}
\begin{singlespace}\noindent
A previsão da volatilidade cambial é relevante para a mensuração de riscos financeiros e a tomada de decisões nos mercados. 
Nesse contexto, este trabalho compara o desempenho preditivo de abordagens model-based e data-driven na previsão da volatilidade 
diária da taxa de câmbio USD/BRL. Para isso, são estimados os modelos GARCH(1,1), EGARCH(1,1) e Random Forest, utilizando dados 
diários da série cambial obtida no Federal Reserve Bank of St. Louis (FRED), referentes ao período de janeiro de 2000 a maio de 2026. 
A volatilidade é representada pelo quadrado dos retornos logarítmicos diários, e as previsões são avaliadas fora da amostra por meio das 
métricas MSE, RMSE, MAE e QLIKE e do teste de Diebold-Mariano. Na janela expansiva, o Random Forest apresentou os menores valores de MSE, 
RMSE e QLIKE, enquanto o EGARCH obteve o menor MAE. Entretanto, não foram identificadas diferenças estatisticamente significativas entre os 
modelos. Na análise de robustez com janela rolante de cinco anos, o GARCH apresentou os melhores resultados em todas as métricas, com 
diferenças significativas em relação ao Random Forest para MSE e QLIKE. Os resultados indicam que a flexibilidade dos métodos de aprendizado 
de máquina não garante ganhos preditivos e que o desempenho relativo das abordagens depende da estratégia de estimação adotada.
\par\vspace{\baselineskip}
\noindent Palavras-chave: volatilidade cambial; previsão de volatilidade; GARCH; Random Forest; aprendizado de máquina.
\end{singlespace}

%% --- ABSTRACT --------------------------------------------------------
\cleardoublepage
\begin{otherlanguage}{english}
\chapter*{Abstract}
\begin{singlespace}\noindent
Predicting exchange rate volatility is relevant for financial risk measurement and decision-making. In this context, this study compares the 
predictive performance of model-based and data-driven approaches in forecasting the daily volatility of the USD/BRL exchange rate. To this end, 
the GARCH(1,1), EGARCH(1,1), and Random Forest models are estimated using daily exchange rate data obtained from the Federal Reserve Bank of St. 
Louis (FRED), covering the period from January 2000 to May 2026. Volatility is represented by squared daily log returns, and the forecasts are 
evaluated out of sample using the MSE, RMSE, MAE, and QLIKE metrics, as well as the Diebold-Mariano test. Under the expanding window, Random Forest 
achieved the lowest MSE, RMSE, and QLIKE values, while EGARCH yielded the lowest MAE. However, no statistically significant differences were found 
among the models. In the five-year rolling window robustness analysis, GARCH achieved the best results across all metrics, with statistically 
significant differences relative to Random Forest for MSE and QLIKE. The results indicate that the flexibility of machine learning methods does not 
guarantee predictive gains and that the relative performance of the approaches depends on the estimation strategy adopted.
\par\vspace{\baselineskip}
\noindent Keywords: exchange rate volatility; volatility forecasting; GARCH; Random Forest; machine learning.
\end{singlespace}
\end{otherlanguage}

%% --- LISTAS (obrigatórias quando houver mais de 5 itens) -------------
\cleardoublepage
\listoffigures
\cleardoublepage
\listoftables

%% --- SUMÁRIO (último elemento pré-textual) ---------------------------
\cleardoublepage
\tableofcontents

%% ======================================================================
%% ELEMENTOS TEXTUAIS (numeração de página começa aqui)
%% ======================================================================
\cleardoublepage
\paginacaonumerada

\subimport{capitulos/}{01_introducao.tex}
\subimport{capitulos/}{02_revisao_literatura.tex}
\subimport{capitulos/}{03_metodologia.tex}
\subimport{capitulos/}{04_base_dados_estrategia.tex}
\subimport{capitulos/}{05_resultados.tex}
\subimport{capitulos/}{06_conclusao.tex}

%% ======================================================================
%% ELEMENTOS PÓS-TEXTUAIS
%% ======================================================================

%% --- REFERÊNCIAS -------------------------------------------------------
\cleardoublepage
\printbibliography[heading=bibintoc,title={Referências}]

%% --- APÊNDICES ---------------------------------------------------------
\iniciarapendices
\subimport{capitulos/}{08_apendice.tex}
\subimport{capitulos/}{09_declaracao_ia.tex}

\end{document}